\chapter{面向五类特征的统一模型设计}
\label{chap:3}
\chapterguide{围绕参数及其作用设计中间结构，保持与现有CADExchange定义一致。可选扩展与尚未由三系统完整读写的字段要单独标注。}

\section{统一模型的设计原则与总体结构}
\label{sec:3-1}

\subsection{基于几何构造与参数调查的设计过程}
\label{subsec:3-1-1}
\outlinehint{按构造分析、系统参数调查、公共含义归纳和读写反馈设计结构，避免把原生接口字段机械取并集。}

\subsection{分层结构与特征历史组织}
\label{subsec:3-1-2}
\outlinehint{划分基础几何、标识引用、草图基准、成形特征和配套证据层；说明有序特征列表与引用的关系，不扩展为已实现通用依赖求解器。}

\section{基础几何、单位与引用模型}
\label{sec:3-2}

\subsection{基础几何与量纲约定}
\label{subsec:3-2-1}
\outlinehint{定义点、方向、局部坐标系及模型单位；说明长度换算、角度解释和无量纲数据的处理规则。}

\subsection{对象标识、引用角色与几何定位信息}
\label{subsec:3-2-2}
\outlinehint{定义引用对象类别、局部标识、作用角色、几何位置和辅助点等信息，分别承担依赖描述与目标选择定位。}

\subsection{特征依赖与几何快照的组织}
\label{subsec:3-2-3}
\outlinehint{区分统一模型中的构造关系与运行阶段的几何快照和源目标对应记录，交代数据产生时机及后续用途。}

\section{草图与基准面的统一结构}
\label{sec:3-3}

\subsection{草图图元、坐标与构造属性}
\label{subsec:3-3-1}
\outlinehint{列出支承关系、坐标系、图元、局部标识和构造属性，逐项解释用途；区分已有具体图元与样条等未完成范围。}

\subsection{草图约束、尺寸与外部引用}
\label{subsec:3-3-2}
\outlinehint{组织约束类型、参与对象、作用位置、数值与外部来源，说明定义完整不代表全部约束已在三系统恢复。}

\subsection{基准面方法、参考与约束参数}
\label{subsec:3-3-3}
\outlinehint{定义建面方法、参考对象、距离、角度、方向和可选实际几何；明确约束组合才是目标构造的输入。}

\section{拉伸与旋转的统一结构}
\label{sec:3-4}

\subsection{材料操作、终止条件与薄壁共享参数}
\label{subsec:3-4-1}
\outlinehint{定义加料与切除、两端终止类型和参考、偏移启用状态及两侧薄壁偏置；对依赖特征上下文的量纲作明确限定。}

\subsection{拉伸轮廓、方向与作用范围}
\label{subsec:3-4-2}
\outlinehint{给出轮廓引用、实际方向、两端范围、对称总量及附属选项的参数作用表。}

\subsection{旋转轮廓、轴线与角度范围}
\label{subsec:3-4-3}
\outlinehint{分别表达轴的位置和来源、公共正向及两侧角度，保留参考终止所需的方向与辅助信息。}

\section{扫掠的统一结构}
\label{sec:3-5}

\subsection{三类截面及圆管参数表达}
\label{subsec:3-5-1}
\outlinehint{分别定义独立草图引用、内嵌草图和参数圆形截面，说明外半径、内半径及各路径的实际使用范围。}

\subsection{路径、有序引用与起止信息}
\label{subsec:3-5-2}
\outlinehint{定义主路径和相关引用、起点终点、分段顺序及闭合标记，说明缺失数据与明确标记的区别。}

\subsection{截面放置与导引控制参数}
\label{subsec:3-5-3}
\outlinehint{组织截面放置、沿路径的方向控制、导引及薄壁参数；对高级控制保留结构范围与生产范围的差异。}

\section{模型存储与合法性检查}
\label{sec:3-6}

\subsection{序列化与可选参数处理}
\label{subsec:3-6-1}
\outlinehint{说明文本化存储、标识一致性，以及未提供、未启用、零值和未知值的区别；不增加未经实现的格式能力。}

\subsection{参数、引用与组合条件检查}
\label{subsec:3-6-2}
\outlinehint{列出尺寸有效性、必要参考、截面与路径条件等检查，并区分公共合法性与目标系统可执行性。}

\section{本章小结}
\label{sec:3-7}

\outlinehint{归纳本章能够支持的认识，并说明与下一章的衔接；不重复列举所有接口或扩展未证明结论。}
